\scanpage{224}
% slanted
\margin{Cela marche dès que l'on a $\deleted{\varphi}: S\underset{\varphi'}{\overset{\varphi}{\rightleftarrows}}S^\vee$
avec $\varphi'\circ\varphi=\mathrm{id}_{\unsure{S^\vee|S}}$, alors
$(\varphi,\varphi')$ est \uline{compatible} avec $R,R'$ ssi $\varphi$ est une jauge}
\[
 \varphi\times\varphi'\text{ transforme }R\text{ en }R'\overset{\mathrm{déf}}={}^sR
 \qquad\text{ssi }\varphi\text{ est une jauge.}\tag*{(8)}
\]
Ce qui donne une correspondance biunivoque \deleted{\unsure{une}}
\[
 \underset{\text{strictes}}{\mathrm{Jauges}}(\Gamma)\simeq
 \begin{aligned}[t]&\text{Ens des isos }(\varphi,\varphi'):\Gamma\longrightarrow\Gamma'\\
 &\text{tels que }\varphi\deleted{\unsure{\varphi}},\varphi'\text{ soient inverses}\\
 &\text{l'un de l'autre, où }\Gamma'\text{ est le}\\
 &\text{prégraphe }(S^\vee,S,{}^sR)\text{ dual de }\Gamma.
 \end{aligned}\tag*{(9)}
\]

Si $\varphi$ est une jauge de $\Gamma$, alors
\[
 \Gamma_\varphi=(S,A_\varphi),\quad\text{où }\ 
 \underset{\substack{\Vert\\\{\{s,t\}\mid\{s,\varphi(t)\}\in R\}}}{A_\varphi}\subset\mathcal P(S)
 \ ,\quad\begin{aligned}[t]&\text{est un graphe}\\&\text{(avec des arêtes}\\
 &\text{qui peuvent être}\\&\text{ponctuelles)}
 \end{aligned}\tag*{(10)}
\]
la relation symétrique associée à cette str. de
graphe étant
\[
 R_\varphi=(\mathrm{id}_S\times\varphi)^{-1}(R)=(\varphi\times\mathrm{id}_S)^{-1}(R)
 =\{(s,t)\in S\times S\mid\{s,\varphi(t)\}\in R\}\tag*{(11)}
\]
-- c'est le \uline{graphe associé à} $\varphi$, \deleted{\unsure{noté} $\Gamma_\varphi$ Si $R$ est sép}
\inserted{On trouve ainsi} \deleted{sép} une application
\[
 \mathrm{Jauges}(\Gamma)\longrightarrow
 \begin{aligned}[t]&\text{ens des structures}\\&\text{de graphes }\gamma\text{ sur }S\\
 &\text{telles que l'ens des étoiles }\mathcal E_\gamma\\
 &\text{correspondants soit }\deleted{\text{égal à}}\inserted{\text{contenu dans}}\\
 &\mathcal E=\mathrm{Im}(S^\vee\longrightarrow\mathcal P(S))\subset\mathcal P(S)
 \end{aligned}\tag*{(12)}
\]
% slanted
\margin{$\mathcal E\inserted{\subset\mathcal P(S)}$ s'appelle l'ens des \uline{étoiles} du prégraphe, on définit de $=$
$\mathcal E^\vee\subset\mathcal P(S^\vee)$, l'ens des coétoiles (une étoile est un ens de sommets, une coétoile un ens. de cosommets\ldots)}
Cette application étant \uline{bijective} si $\Gamma$ est séparée en $S^\vee$,
i.e. si on peut identifier $S^\vee$ à \deleted{\unsure{l'}} un sous-ens de $\mathcal P(S)$
\[
 S^\vee\hookrightarrow\mathcal P(S).\tag*{(13)}
\]
\deleted{Dans ce cas, \unsure{on}} On trouve une \unsure{réflection|restriction|application} \unsure{en dessous|évidente|en deçà}
\[
 \mathrm{Jauges\ surj}(\Gamma)\longrightarrow
 \begin{aligned}[t]&\{\text{ens des str. de}\\&\text{graphes }\inserted{\gamma}\text{ sur }S\\
 &\text{telles que }\mathcal E_\gamma=\mathcal E\\
 &[\text{i.e. ayant un ens}\\&\text{d'étoiles donné}
 \end{aligned}\tag*{(14)}
\]
On trouve aussi
% archive sheet number: 223

